% Folder 93, batch 04 (archivists' pages 61-80).
% Pass: Opus 5.5 (claude-opus-5-5), 2026-10-10 — first pass, unchecked against the pages by a human.
%
% Pages 62, 64 and 66 are unrelated typescript versos (an administrative
% report on university research), 68-69 a list of names and an address,
% 70 and 78 blank; all are absent. Pages 61-65 continue an argument on Euler
% characteristics of locally conic spaces begun before this batch; pages
% 71-80 are a separate run on differential invariants (Schwarzian
% derivative) of étale maps of curves.
\documentclass[11pt,a4paper]{article}
\input{../preamble/grothendieck}

\folder{93}
\batch{4}
\pages{61}{80}
\foldertitle{Connexions de Gauss-Manin et équations de Picard-Fuchs. Opération de Cartier : copies de tapuscrit annoté (s.d.), tiré à part (1981), notes manuscrites (s.d.).}
\dating{[1972]-1981}
\watermark{Édition de démonstration}

\begin{document}

\page{61}
\note{la page s'ouvre au milieu d'un argument commencé avant ce lot : les classes $\mathcal{E}_\chi$, $\mathcal{E}_{\chi_!}$ et les conditions $(\psi_1)$, $(\psi_2)$ y sont supposées connues}
\uncertain{Cela mène à la} question si \struck{$X$} $Y$ \add{(cpct)} \add{$(\psi_1)$} admet des recouvrements arbitrairement fins \add{finis} par des $Y_i$ \struck{tels} \add{\uncertain{fermés}} tels que les intersections $Y_{i_1} \cap Y_{i_2} \cap \dots \cap Y_{i_r}$ $(r \geq 1)$ de tels $Y_i$ soient $\in \mathcal{E}_{\chi_!}$ …

Dans le cas des $\chi$, \uncertain{on utilise le seul symbole} d'\emph{homologie}, \uncertain{on est} ramené :

\emph{Lemme} Soit $R$ syst.\ local \ill{} \uncertain{homotopique} sur $Y \in \mathcal{E}_\chi$, \add{finis}, de rang $b$. Alors
\[
\chi(Y, R) \overset{\mathrm{def}}{=} \sum_p (-1)^p H_p^{\flat}(Y, R)
\]
\note{sous $H_p^{\flat}(Y,R)$, une accolade : « le \uncertain{même} que $H^p(Y, \check{R})$ » ; le signe en exposant, lu $\flat$, est incertain}
existe et
\[
\chi(Y, R) = b\,\chi(Y)
\]

Cela montre pourquoi \ill{} faire l'hyp.\ \struck{\ill{}} $Y \sim Y'$ satisfaisant la \emph{condition} $(\psi_2)$.

\struck{On va se borner}

ⓗ Définissons $\mathcal{C}_{-1} =$ ens.\ réduit à l'espace vide

et si $n \geq 0$, $\mathcal{C}_n =$ \struck{catégorie} \add{ens.}\ des espaces cpts $X$ tels que $\forall x \in X$ il y ait un voisinage $C_x$ tel qu'il existe un espace compact $S \in \mathcal{E}_{n-1}$ et un homéom.\ $\mathrm{Cône}(S) \simeq C_x$ transformant le sommet du cône en $x$.

\marginal{NB $\mathcal{E}_0$ espaces discrets ; $\mathcal{E}_1$ 1-complexes topologiques ; $\mathcal{E}_2$ espaces loc.\ triangulables de dim.\ $\leq 2$ [ils sont triangulables dans le cas cpt] ; $\mathcal{E}_3$ espaces loc.\ triangulables de dim.\ $\leq 3$ [\ill{} triangulables dans le cas cpt] ; $\mathcal{E}_4$ \uncertain{espaces loc.\ triangulables} de dim.\ $\leq 4$ (Moïse) ; s'ils sont cpts, il n'est pas clair qu'ils soient triangulables …}

$\mathcal{C}_\infty = \bigcup \mathcal{C}_n$ (espaces loc.\ coniques)

$\mathcal{E}_\infty =$ \struck{\ill{}} des espaces loc.\ cpts \add{séparés} $X$ tels que $\hat{X}$ soit loc.\ conique i.e.\ soit $\in \mathcal{C}_\infty$.

\emph{Prop} $\mathcal{C}_\infty \subset \mathcal{E}_\chi \cap \mathcal{E}_{\chi_!}$

Soit $X \in \mathcal{E}_\infty$. Si $X$ est cpt, \add{donc $X \in \mathcal{C}_\infty$} comme $X$ est loc.\ contractile, on sait que $X \in \mathcal{E}_\chi$, donc $X \in \mathcal{E}_{\chi_!}$. Cas \struck{$X$ \ill{}} \add{général} : $X$ loc.\ cpt $= \hat{X} - (\omega)$, avec $\hat{X} \in \mathcal{C}_\infty$, on a $\hat{X} \in \mathcal{E}_{\chi_!}$, donc $\hat{X} - \omega = X \in \mathcal{E}_{\chi_!}$ par (b). Pour prouver $X \in \mathcal{E}_\chi$, prenons un $C_x \simeq C(S)$, prenons un homothétique $\frac{1}{2} C_x = T$, et

\page{63}
\note{suite de la démonstration de la p.~61, la p.~62 étant sans rapport}
Soit $X' = \hat{X} - \mathring{T}$, qui est cpct. \struck{$X \in \mathcal{C}$.} On vérifie
\marginal{vérifier}

a) $X' \in \mathcal{C}_\infty$ donc $X' \in \mathcal{E}_\chi$

b) $X' \sim X = \hat{X} - \omega$ donc $X \in \mathcal{E}_\chi$,

\emph{NB} On trouve
\[
\chi(X) = \chi(\hat{X} - \mathring{T}) = \chi(\hat{X}) - \chi(T) + \chi(\partial T)
\]
\note{sous $\chi(T)$, un renvoi biffé et peu lisible}
et $\chi(T) = 1$ car $T$ contractile, donc
\[
\chi(X) = \chi(\hat{X}) - 1 + \chi(\partial T) = \chi_!(X) + \chi(\partial T), \qquad \partial T = S_\omega,
\]
i.e.
\[
\boxed{\chi(X) - \chi_!(X) = \chi(S_\omega)}
\]
où $\chi(S_\omega)$ « exprime la "structure locale de $X$ à l'infini" ».

\page{65}
\textbf{Th} Sur la cat.\ $\mathcal{E}_\infty$ des espaces loc.\ cpts séparés et \struck{\ill{}} « loc.\ coniques partout » (y compris à l'$\infty$), on peut trouver deux fonctions uniques $\chi_!$, $\chi$ (à valeurs $\in \mathbb{Z}$) satisfaisant les conditions (où tous les espaces écrits sont supposés $\in \mathcal{E}_\infty$) :

ⓐ $\chi_!(X) = \chi(X)$ si $X$ cpt

ⓑ $Y$ fermé dans $X$ $\Longrightarrow$ $\chi_!(X) = \chi_!(Y) + \chi_!(X - Y)$

ⓒ $X'$, $X''$ fermés dans $X$ $\Longrightarrow$ $\chi(X) = \chi(X') + \chi(X'') - \chi(X' \cap X'')$\note{la condition $X = X' \cup X''$ n'est pas écrite}

ⓓ $\begin{cases} X \sim X' \Longrightarrow \chi(X) = \chi(X') \\ X \simeq X' \longrightarrow \chi_!(X) = \chi_!(X') \end{cases}$

ⓔ $\chi(\mathrm{pt}) = 1$

ⓕ $\chi_!(X \times Y) = \chi_!(X)\,\chi_!(Y)$

$\chi(X \times Y) = \chi(X)\,\chi(Y)$

ⓖ si $f : X \to Y$ fibration loc.\ triviale de fibre $F$,
\[
\begin{cases} \chi_!(X) = \chi_!(Y)\,\chi_!(F) \\ \chi(X) = \chi(Y)\,\chi(F) \end{cases}
\]
\marginal{j'ignore si c'est vrai sans telle restriction}
\uncertain{vrai} au moins si $Y$ est un polyèdre relatif (diff.\ de deux polyèdres, l'un plongé dans l'autre linéairement…)

Si $X \in \mathcal{E}_\infty$ et tel que $\hat{X} \in \mathcal{E}_\infty$, …

ⓗ
\[
\begin{cases} \chi_!(X) = \chi(\hat{X}) - 1 \\ \chi(X) - \chi_!(X) = \chi(S) \end{cases}
\]
où $S$ est la « frontière à l'infini » de $X$.

Enfin, la restriction de \struck{$\mathcal{E}$} $\chi$ et $\chi_!$ aux polyèdres relatifs sont uniquement déterminées par les propriétés ⓐ à ⓔ, et on a, si $X$ muni d'une filtration \ill{},
\[
\chi_!(X) = \sum (-1)^i f_i \qquad \text{(formule d'Euler)}
\]
\marginal{j'ignore si $\chi$ et $\chi_!$ sont caractérisées par a) à e) sur tout $\mathcal{E}_\infty$}

\page{67}
\note{liste de thèmes, seule sur le feuillet ; une accolade réunit les trois lignes sur les revêtements}

orientations

\uncertain{plongements}

$H^i$, $\chi$

rev.\ et $\pi_1$ ;
rel.\ entre $\pi_1$ et th.\ de Galois ;
rev.\ et fonctions de var.\ complexes

Conj.\ de Weil

\page{71}
\note{changement de sujet et de papier (feuillets perforés, à l'encre bleu-noir) : invariants différentiels d'une courbe analytique ; les marges de gauche sont couvertes d'ajouts écrits en long, partiellement lus}

\textbf{Th} Ⓘ Soit $X$ une courbe analytique \emph{régulière}, $\mathcal{E}_X$ $(\mathcal{P}_X)$ le faisceau des applications étales d'ouverts de $X$ dans $\mathbb{E}^1(\mathbb{C})$ $(\mathbb{P}^1(\mathbb{C}))$. Il existe un morphisme
\[
G : \mathcal{E}_X \times \mathcal{E}_X \longrightarrow \underline{\omega}_{X/\mathbb{C}}^{\otimes 2} \qquad (\underline{\omega}_{X/\mathbb{C}} = \underline{\Omega}^1_{X/\mathbb{C}})
\]
ayant les propriétés suivantes :

a) $G$ est un op.\ différentiel (non linéaire), i.e.\ en termes d'une uniformisante locale $t$ \add{(définie sur un ouvert $U$)},
\[
G(f,g) = R(f, f', \dots, f^{(n)}, g, g', \dots, g^{(n)})\, dt^2
\]
où $R$ est une fonction rationnelle en $2n$ variables, à coefficients dans $\Gamma(U, \mathcal{O}_U)$. De plus, \add{si $R = P/Q$, $P$ et $Q$ polynômes,} \struck{\ill{}} \struck{\ill{} \ill{}} ce qui implique que $Q(f, f', \dots, f^{(n)}; g, g', \dots, g^{(n)})$ n'est pas identiquement nulle, \struck{\ill{} \ill{}} \add{ce qui signifie que $Q$ ne s'annule identiquement} pour aucunes $f$ et $g$ \uncertain{ouvertes}, \uncertain{non} loc.\ constantes sur l'une ou l'autre, et que $P(f, \dots, g, \dots)/Q(f, \dots, g, \dots)$ n'a pas de pôles \uncertain{quand} $f'$, $g'$ ont des zéros.

[On doit pouvoir montrer que, si on prend \ill{} \add{premiers entre eux}, \ill{} \ill{}, $Q = f'^{\,?} g'^{\,?}$, donc
\[
G(f,g) = \frac{1}{f'^{\,?} g'^{\,?}}\, P(f, f', \dots, f^{(n)}; g, g', \dots, g^{(n)}) \;]
\]
\note{les exposants de $f'$ et $g'$ sont des signes illisibles}

b) $G(f,g) + G(g,h) = G(f,h)$

c) \struck{\ill{}} \struck{si $a \in \mathbb{C}^*$, $b \in \mathbb{C}$} \struck{et}
\[
\begin{array}{ll}
G(f, ag) = G(f,g) & a \in \mathbb{C}^* \\
G(f, g+b) = G(f,g) & b \in \mathbb{C} \\
G(f, 1/g) = G(f,g) & \text{si } g \text{ sans zéros}
\end{array}
\]

d) \struck{\ill{}} Si $G(f,g) = 0$ et $f$, $g$ coïncident à l'ordre 2 en un point $x \in U$ ($f, g \in \Gamma(U, \underline{\mathcal{O}}_U)$) alors $f = g$ si $U$ connexe.

\marginal{ⓔ si $f \in \Gamma(U, \mathcal{E}_U)$, $g \mapsto G(f,g)$ de $\mathcal{E}_U$ dans $\underline{\omega}_U^{\otimes 2}$ est un épi}

Deux tels morphismes $G'$ et $G$ sont liés par
\[
G' = \varphi\, G, \quad \text{avec } \varphi \in \Gamma(X, \underline{\mathcal{O}}_X^*)
\]
\marginal{Ici, au lieu des conditions d), e), il suffit de \ill{} d)}

Ⓘⓘ Si, au lieu de prendre $X$ fixe, on prend $X$ variable, avec \uncertain{les} morphismes étales des $X$, on a \ill{} résultat, et l'indétermination sera
\[
G' = c\, G \qquad c \in \mathbb{C}^*
\]

\marginal{NB c) implique que $G$ se factorise en $\mathcal{P}_X \times \mathcal{P}_X \to \underline{\omega}_X^{\otimes 2}$, ou $H \backslash \mathcal{P}_X \times_H \mathcal{P}_X \to \underline{\omega}_X^{\otimes 2}$, \ill{} ; b) donne la \uncertain{relation} $H \backslash \mathcal{P}_X \to \mathcal{T}$ \ill{} ; $H = GP(1, \mathbb{C})$. \ill{} $\mathcal{P}_X \to \mathcal{T}_X$ \ill{}}
\note{les autres ajouts de la marge gauche, écrits en long et en partie biffés, n'ont pas pu être lus}

\page{72}
[Ici, on montre que \add{la condition a) implique que} pour tout $X$ fixé, l'expression de a), $G(f,g) = R(f, f', \dots, f^{(n)}; g, g', \dots, g^{(n)})$, doit avoir des coefficients $\in \mathbb{C}$, indépendants de $X$, donc
\[
R \in \mathbb{C}(F_0, F_1, \dots, F_n; G_0, \dots, G_n)
\]
et en fait,
\[
R(\ ) = \underbrace{P(F_0, \dots; G_0, \dots)}_{\in\, \mathbb{C}[F_0, \dots, G_n]} / F_1^{\,?}\, G_1^{\,?} \;]
\]

On peut prendre
\[
G(f,g) = \left( \left[ -2\, \frac{g'''}{g'} + 3 \left( \frac{g''}{g'} \right)^2 \right] - \left[ -2\, \frac{f'''}{f'} + 3 \left( \frac{f''}{f'} \right)^2 \right] \right) dt^2
\]

De plus, on trouve un torseur canonique $\mathcal{T}_X$ \struck{\ill{}} \uncertain{sous} $\underline{\omega}_{X/\mathbb{C}}^{\otimes 2}$, dépendant fonctoriellement de $X$ (pour morphismes étales), et \struck{\ill{}}
\[
c^X : \mathcal{P}_X \longrightarrow \mathcal{T}_X
\]
tel que
\[
\begin{cases}
G(f,g) = c^X(g) - c^X(f) \\
c^X \text{ se factorise en } {}_H\backslash \mathcal{P}_X \to \mathcal{T}_X \\
\text{ce dernier morphisme est un iso}
\end{cases}
\]

III) On prend maintenant $X$ \add{\uncertain{schéma}} (\uncertain{variété} \struck{lisse} \uncertain{lisse} de \uncertain{dim.\ rel.\ égale à} 1) sur des schémas variables, et on cherche pour tout tel $X$
\[
\mathcal{E}_X \times \mathcal{E}_X \xrightarrow{\ G\ } \omega^{2}_{X/S}
\]
de façon compatible avec chgt de base $S' \to S$, aux morphismes étales $X' \to X$ \uncertain{au-dessus} de $S$, et satisfaisant les conditions qui correspondent à a) b) c) (\struck{si} $S$ \uncertain{courbe} de car.\ 0, et e) \ill{}), plus les conditions d) \struck{\ill{} \ill{}} \struck{« analytique complexe »}.
\marginal{la forme \uncertain{précédente} \ill{} \ill{} $G(f,g)$ \ill{} \uncertain{cont.} $f'$ \ill{}}
[Mais quand on parle « d'op.\ différentiel », il faut maintenant remplacer un polynôme $R$ en $f, f', f'', \dots, g, g', g'', g'''$ par un polynôme en
\[
f,\ f' = D_t^{(1)} f,\ D_t^{(2)}(f), \dots, D_t^{(n)} f,\ g,\ g' = D_t^{(1)} g,\ D_t^{(2)} g, \dots, D_t^{(n)} g.
\]
\struck{On voit \ill{}} A priori, les coeff.\ sont dans $\Gamma(U, \underline{\mathcal{O}}_U)$ (où $t$ est \uncertain{uniformisante} définie sur $U$) mais si on \struck{\ill{}} \uncertain{les} fait varier $X$, on voit que les coefficients doivent être dans $\mathbb{Z}$ (ou dans $\Gamma(S_0, \underline{\mathcal{O}}_{S_0})$ si on a fixé un schéma de base commun $S_0$ pour les $S, X$ …). On trouve encore l'existence d'un tel, et l'unicité à un facteur $\varepsilon \in \mathbb{Z}^*$ près, i.e.\ $\varepsilon = \pm 1$ : \emph{unicité au signe près} ! Ici, on ne peut pas prendre l'expression $G(f,g)$ précédente, car elle a des coeff.\ 2 et 3 ! Il faut prendre $\frac{1}{12} G(f,g)$, savoir
\[
G_0(f,g) = \left[ \left( -\frac{D_t^{(3)} g}{D_t^{(1)} g} + \left( \frac{D_t^{(2)} g}{D_t^{(1)} g} \right)^2 \right) - \left( -\frac{D_t^{(3)} f}{D_t^{(1)} f} + \left( \frac{D_t^{(2)} f}{D_t^{(1)} f} \right)^2 \right) \right] dt^2
\]

Pour déterminer ce choix (parmi les deux possibles au signe près) on va calculer p.ex.
\[
G_0(t, t^2) = \text{\struck{$\frac{1}{2t}$}}\ \frac{1}{2t}\, dt^2 \qquad \left(\text{au lieu de } -\frac{1}{2t}\, dt^2\right)
\]
\note{le calcul direct donne $\frac{1}{4t^2}\,dt^2$ ; on transcrit tel qu'écrit}
[Je ne fais pas d'affirmation sur ce dénominateur, \uncertain{parce que} 2 n'est pas inv.\ sur $X$ : alors $t^2$ n'est jamais \uncertain{une} uniformisante !] On \uncertain{trouve} \uncertain{aussi} quelque chose qui ait un sens aussi en car.\ 2 :
\[
G_0(t, t^3) = \underbrace{\left( -\frac{1}{3t^2} + \left( \frac{3t}{3t^2} \right)^2 \right)}_{\frac{1}{t^2}\left(-\frac{1}{3} + 1\right)} dt^2 = \text{\struck{$\frac{1}{3t^2}$}}\ \frac{2}{3}\, \frac{1}{t^2}\, dt^2
\]
[mais c'est inutile, car en car.\ 2 les \struck{\ill{}} \add{deux choix} possibles sur $\mathbb{Z}$ coïncident !].

\emph{Remq} 1) On peut se demander quel est le rôle joué ici par les \uncertain{bijections}, et s'il n'y aurait pas moyen d'éliminer ce rôle. On est alors conduit, dans I), à se donner un faisceau analytique cohérent $\mathcal{F}$ sur $X$, et à supposer $G' : \mathcal{E}_X \times \mathcal{E}_X \to \mathcal{F}$. On trouve alors qu'un tel $G'$ s'écrit de façon unique (si $G$ est connu dès l'\uncertain{énoncé} initial, p.ex.\ celui de II)) comme $G' = uG$, où $u : \underline{\omega}_X^{\otimes 2} \to \mathcal{F}$, et \uncertain{on voit} \uncertain{que} $u$ est un \emph{iso} ! Donc $\underline{\omega}_X^{\otimes 2}$ joue ici un rôle \emph{universel}. Dans II) on pourrait se donner un $\mathcal{F}_X$ dépendant

\page{73}
fonctoriellement de $X$, mais on trouverait encore un iso fonctoriel en $X$, $\mathcal{F}_X \simeq \underline{\omega}^{\otimes 2}_{X/\mathbb{C}}$. \uncertain{De même} dans (III). Ainsi les 2 formes \uncertain{sont} \uncertain{essentielles}.

2) \struck{$G$} Le torseur $\mathcal{T}$ dont il est question dans II (\uncertain{et en marge} du I) dépend du choix de $G$ (\uncertain{tel} de $\mathcal{F}$, \uncertain{si on veut} — c'est donc un torseur \uncertain{sous} $\mathcal{F}$) mais les \add{torseurs associés aux} différents choix sont \uncertain{canoniquement} \ill{}, de façon précise
\[
\mathcal{T}^G \xrightarrow{\ \sim\ } \mathcal{T}^{\varphi G}
\quad \text{au-dessus de} \quad
\underline{\omega}_X^{\otimes 2} \xrightarrow{\ \varphi \cdot \mathrm{id}\ } \underline{\omega}_X^{\otimes 2}
\]
\note{les deux flèches sont écrites l'une au-dessus de l'autre, reliées par un trait vertical}
i.e.\ $\mathcal{T}^{\varphi G}$ est « multiple » de $\mathcal{T}^G$ par $\varphi$. \struck{A} Ainsi, le choix « le plus fin » de $\mathcal{T}$ est $\mathcal{T}^{G_0}$ (où $G_0$ est défini en (\struck{II}) III), alors le $\mathcal{T}^G$ correspondant au $G$ quelconque de (III) est un multiple (12 \uncertain{inverse} \ill{}) de $\mathcal{T}^{G_0}$.

3) On \uncertain{pourrait} se \struck{demander} \add{poser} la question de caractériser $G_0(f,g)$ en une \emph{caractéristique donnée} $p > 0$. On \uncertain{pourra} prendre les conditions a) b) c), plus d') ($G_0(f,g)$ n'est pas id.\ nulle).
\marginal{\uncertain{faux} \ill{} \ill{} \ill{} d'ordre $p$ ?}
\emph{\ill{}} caractérise $G_0(f,g)$ \emph{à une multiplication près} par une constante $c \in \mathbb{F}_p^*$. \ill{} \ill{} si on fixe un schéma de base absolu (\uncertain{pour tous les} $S$ et les $X$), l'indétermination \ill{} dans $\Gamma(S, \underline{\mathcal{O}}_S^*)$. Si $S$ n'est \emph{\uncertain{intègre}}, il faut supposer que, en chaque \ill{}, \ill{} des $X$ \uncertain{lisses} \ill{} (\uncertain{rel.\ à sa base}), \ill{}
\marginal{d'ordre \ill{} identiquement nul \ill{}}

\page{74}
Les développements précédents étaient centrés sur le \emph{groupe} $H = GP(1, \mathbb{C})$ qui opère sur $\mathbb{P}^1(\mathbb{C})$ (ou $GP(1)_{\mathbb{Z}}$ sur $\mathbb{P}^1_{\mathbb{Z}}$ …), et les invariants différentiels, sous l'action de ce \emph{groupe}, des fonctions analytiques complexes (\emph{cas} \uncertain{modèle} — \uncertain{vaut} aussi dans l'analytique réel, en les ($C^\infty$, ou $C^3$)). On va passer en revue les développements analogues pour d'\emph{autres groupes}.

a) \emph{Groupe des translations}, \emph{opérant sur} $\mathbb{E}^1$.

On a l'opération (dans l'analytique complexe p.ex.)
\[
d : \underline{\mathcal{O}}_X \longrightarrow \underline{\Omega}^1_{X/\mathbb{C}} = \underline{\omega}_X
\]
épi \struck{(X lisse de dim.\ 1) et \ill{}} \add{noyau exactement} sur les fonctions loc.\ constantes, donc on prend
\[
A(f,g) = dg - df, \qquad A : \underline{\mathcal{O}}_X \times \underline{\mathcal{O}}_X \longrightarrow \underline{\omega}_X
\]
\marginal{Ici, $\underline{\mathcal{O}}_X$ mauvais, vaut mieux prendre $\mathcal{E}_X$ plutôt !}
On trouve l'analogue de a) b) dessus, plus
\[
A(f, g+c) = A(f,g) \qquad c \in \mathbb{C}
\]
plus
\[
A(f,g) = A(f,h),\ g(x) = h(x) \Longrightarrow g = h \text{ au vois.\ de } x
\]
et
\[
g \longmapsto A(f,g) \qquad \underline{\mathcal{O}}_X \to \underline{\omega}_X \text{ épi}
\]
(car $X$ lisse de dim.\ 1 \uncertain{par définition}) — \add{(si non, il est trivial)}

\uncertain{Ici}, pas la peine d'introduire de torseur, \uncertain{car} $\omega = A(f,g)$, \uncertain{on a} l'expression \ill{} \ill{} \ill{} qui en \uncertain{sont} \ill{}

\page{75}
Si on s'intéresse à $\mathcal{E}_X$ au lieu de $\underline{\mathcal{O}}_X$, il faut remplacer $\underline{\omega}_X$ par $\underline{\omega}_X^*$ (sous-faisceau des bases de $\underline{\omega}_X$, i.e.\ des sections sans zéros …), qui est un torseur \struck{sous} \add{sous} $\underline{\mathcal{O}}_X^*$, et $dg - df$ par \struck{$dg$} $A'(f,g) = \frac{dg}{df} \in \Gamma(X, \underline{\mathcal{O}}_X^*)$.

b) \emph{Groupe opérant sur} $\mathbb{E}^1$, \emph{ou} $\mathbb{E}^{1*} \simeq \mathbb{G}_m$ par \struck{translations} \add{homothéties}. Plus généralement, considérons les \uncertain{invariants} sur un faisceau inversible quelconque $L$, soit $L^*$ le sous-faisceau des bases \add{(torseur sous $\underline{\mathcal{O}}_X^*$)}, on s'intéresse à $L^*/\mathbb{C}^*$. Or une section de $L^*$ est \uncertain{associée à} une connexion sur $L$ : on trouve ainsi un homom.
\[
L^* \xrightarrow{\ D\ } \underline{\mathrm{Conn}}_{X/\mathbb{C}}(L) \qquad (f \longmapsto D_f)
\]
où le dernier membre est un torseur sous $\underline{\omega}_{X/\mathbb{C}}$, et ceci induit un \struck{\uncertain{torseur}} iso
\[
L^*/\mathbb{C}^* \xrightarrow{\ \sim\ } \underline{\mathrm{Conn}}_{X/\mathbb{C}}(L)
\]
Si on ignorait la notion de connexion, \struck{et} on définirait directement $B(f,g)$ $(f, g \in \Gamma L^*)$, par \struck{\uncertain{conn.}} $D_g - D_f$, par
\[
B(f,g) = \frac{d\, g/f}{g/f} \in \Gamma(\underline{\omega}_{X/\mathbb{C}})
\]
avec des propriétés du type a) b) c) d) e). Dans le cas $L = \underline{\omega}_X$, $\underline{\mathrm{Conn}}(L) \simeq \underline{\omega}_{X/\mathbb{C}}$ \uncertain{canon.}, et il suffit de considérer
\[
g \longmapsto \frac{dg}{g} \qquad \underline{\mathcal{O}}_X^* \longrightarrow \underline{\omega}_{X/\mathbb{C}}
\]
qui induit $\underline{\mathcal{O}}_X^*/\mathbb{C}^* \xrightarrow{\ \sim\ } \underline{\omega}_{X/\mathbb{C}}$.
\note{la dernière phrase identifie $L = \underline{\omega}_X$ mais écrit l'exemple avec $\underline{\mathcal{O}}_X^*$, c'est-à-dire $L = \underline{\mathcal{O}}_X$ ; tel quel sur la page}

\page{76}
Cela a mené à $B'(f,g) = \frac{dg}{g} - \frac{df}{f}$ … Si on veut par contre un invariant sur $\mathcal{E}_X^*$, on prendra
\[
g \longmapsto \frac{g}{dg} \qquad \mathcal{E}_X \longrightarrow \underline{\omega}_{X/\mathbb{C}}^{\otimes(-1)}
\]
On \struck{trouve} définit
\[
C^!(f,g) = \text{\struck{$\frac{dg}{g} - \frac{df}{f}$}}\ g(dg)^{-1} - f(df)^{-1}
\]
satisfaisant à a) b) c) [$C^!(f,g) = C^!(f, ag)$ si $a \in \mathbb{C}^*$] d) ($C^!(f,g) = C^!(f,h)$ et $g(x) = h(x)$ en un pt qui \uncertain{n'est pas} un zéro de $g \cdot h$ $\Longrightarrow g = h$) et \struck{e) $[g \mapsto C^!(f,g)]$} pas tout à fait e), car $g(dg)^{-1} - f(df)^{-1}$ (p.ex.\ $g(dg)^{-1}$) \uncertain{ne} peut être id.\ nulle si $g \in \Gamma \mathcal{E}_X$ ! Donc il faut se borner à
\[
\mathcal{E}_X^* = \underline{\mathrm{Et}}_{\mathrm{an}}(X, \mathbb{C}^*)
\]
\add{au quel cas on s'envoie} dans $\underline{\omega}_{X/\mathbb{C}}^{\otimes(-1)\,*}$ par $g(dg)^{-1}$, ou au choix dans $\underline{\omega}_{X/\mathbb{C}}^{*}$ par $dg/g$, et il y a lieu de considérer $\underline{\omega}^*_{X/\mathbb{C}}$ comme un torseur sous $\underline{\mathcal{O}}^*_{X}$. Donc ici on prend
\[
C^{\natural}(f,g) = \frac{dg}{g} : \frac{df}{f} = \frac{dg}{df}\, \frac{f}{g} \in \Gamma(X, \underline{\mathcal{O}}_X^*)
\]
\note{l'indice sur $C$, lu $\natural$, est un signe incertain}

\emph{Dans tous ces cas}, l'espace des valeurs pour les $f, g$ ($\mathbb{P}^1(\mathbb{C})$, $\mathbb{E}^1(\mathbb{C})$, $\mathbb{G}_m(\mathbb{C})$) \struck{étant un espace homogène} sous le groupe d'op.\ envisagé ($GP(1)$, $\mathbb{G}_a$, $\mathbb{G}_m$) ce qui semble condition de succès ($\mathbb{C}^*$ opérant sur $\mathbb{C}$ : \emph{mauvais} !)

c) \emph{Groupe} $\mathrm{Aff}(1) \simeq \mathbb{G}_m \cdot \mathbb{G}_a$ \emph{opérant sur} $\mathbb{E}^1$.

On considère \uncertain{alors}

\page{77}
\begin{tikzcd}[column sep=small]
\mathcal{E}_X \arrow[rr] \arrow[dr] & & \underline{\mathrm{Conn}}(\underline{\omega}_X) \\
& \underline{\omega}^*_{X/\mathbb{C}} \arrow[ur] &
\end{tikzcd}

par $f \longmapsto D_{df}$ (= connexion de $\underline{\omega}_X$ pour laquelle \add{la base} $df$ est horizontale)

on trouve
\[
{}_{\mathrm{Aff}(1,\mathbb{C})}\backslash \mathcal{E}_X \xrightarrow{\ \sim\ } \underline{\mathrm{Conn}}(\underline{\omega}_X)
\]
car (successivement par a))
\[
{}_{\mathbb{G}_a(\mathbb{C})}\backslash \mathcal{E}_X \xrightarrow[\ \sim\ ]{d} \underline{\omega}_X^*
\]
et par b)
\[
{}_{\mathbb{G}_m(\mathbb{C})}\backslash \underline{\omega}_X^* \xrightarrow{\ \sim\ } \underline{\mathrm{Conn}}(\underline{\omega}_X)
\]
On trouve donc
\[
D(f,g) = \frac{d\,\frac{dg}{df}}{\frac{dg}{df}} \in \Gamma\, \underline{\omega}_X
\]
si on a uniformisante $t$, cela fait
\[
D(f,g) = \frac{d\,\frac{g'}{f'}}{\frac{g'}{f'}} = \frac{dg'}{g'} - \frac{df'}{f'} = \left( \frac{g''}{g'} - \frac{f''}{f'} \right) dt
\]
\struck{\uncertain{Remarquons}}

d) \emph{Groupe} \struck{$\mathcal{E}_X$} \add{$\mathbb{G}_m$} « \emph{difféo} » de $\mathbb{G}_m$, $\mathrm{Aut}(\mathbb{G}_m) = \mathbb{Z}/2\mathbb{Z}$, \emph{opérant sur} $\mathbb{G}_m$ ?
\[
f \longmapsto \frac{1}{f} \quad \text{et} \quad f \longmapsto cf \quad (c \in \mathbb{C}^*)
\]
$f$ dans $\underline{\mathrm{Et}}(X, \mathbb{C}^*) = \mathcal{E}_X^*$. On prend
\[
\mathcal{E}_X^* \longrightarrow (\underline{\omega}_X^{\otimes 2})^*, \qquad
f \longmapsto \frac{df}{f} \in \underline{\omega}_X^*, \quad \omega \longmapsto \omega^2
\]
\note{la flèche du bas passe par $\underline{\omega}_X^*$ : $f \mapsto df/f$, puis $\omega \mapsto \omega^2$}

\page{79}
\note{feuillet de calculs, sans phrase d'introduction ; on suit l'ordre de la page}
$\omega$ base de $\underline{\omega}^{\otimes n}$ ($n \in \mathbb{Z}$, \uncertain{$n \neq 0$})

$\omega^{1/n}$ base de $\underline{\omega}$ définie loc.\ à une mult.\ par $\varepsilon$, $\varepsilon^n = 1$ près

$\int \omega^{1/n}$ bien défini mod.\ les transformations $f \mapsto \varepsilon f + b$ ($\varepsilon^n = 1$), donc ${}^{c}\!\!\int \omega^{1/n}$ bien défini.

Si $t$ uniformisante et $\omega = f(t)\, dt^n$, on aura $\omega^{1/n} = f^{1/n}\, dt$ et \struck{\ill{}}
\[
{}^{c}\!\!\int \omega^{1/n} = {}^{c}\!\!\int f^{1/n}\, dt \quad \text{\struck{$= c_t +$}}
\]
\[
c_n(\omega) = c_t + \left[ -\frac{2}{n}\, \frac{f''}{f} + \frac{2n+1}{n^2} \left( \frac{f'}{f} \right)^2 \right] dt^2 \qquad \text{si } n \neq 0
\]

\struck{$G_n(\omega)$}
\[
\begin{aligned}
G_{m,n}(\omega_m, \varpi_n) &= c_n(\varpi_n) - c_m(\omega_m) \\
&= \left[ \left( -\frac{2}{n}\, \frac{g''}{g} + \frac{2n+1}{n^2} \left( \frac{g'}{g} \right)^2 \right) \right. \\
&\qquad \left. - \left( -\frac{2}{m}\, \frac{f''}{f} + \frac{2m+1}{m^2} \left( \frac{g'}{g} \right)^2 \right) \right] dt^2
\end{aligned}
\]
où $\omega_m = f(dt)^m$, $\varpi_n = g(dt)^n$.
\note{dans la seconde parenthèse, la page écrit $(g'/g)^2$ où l'on attendrait $(f'/f)^2$ ; transcrit tel quel. Le coefficient lu $\frac{2n+1}{n^2}$ est d'une lecture incertaine}
\[
\begin{cases}
G_{m,n}(\omega_m, \omega'_n) + G_{n,p}(\omega'_n, \omega''_p) = G_{m,p}(\omega_m, \omega''_p) \\
G_{m,n}(\omega_m, \omega_n) = G_{m,nr}(\omega_m, (\omega_n)^r) \quad \text{si } n \neq 0,\ r \neq 0 \\
G_{m,n}(\omega_m, ds) = G_{m,0}(\omega_m, s) \\
G_{0,0}(t, s) = G(t, s)
\end{cases}
\]
\note{une flèche descend de ces relations vers les deux suivantes}
\[
G_{m,n}(\omega_m, (ds)^n) = G_{m,0}(\omega_m, s) \qquad c'_n(\omega) = n^2\, \frac{c_n(\omega)}{12} \in \mathcal{T}_X^{\otimes n^2}
\]
\note{le dénominateur « 12 » est surchargé, de lecture incertaine ; l'exposant de $\mathcal{T}_X$ est incertain}
\[
G_{m,n}((dt)^m, (ds)^n) = G_{0,0}(t, s)
\]

\struck{$c'_n(\omega)$} ; $m, n$ ; $\delta = m\alpha = n\beta$, $\delta^2 = m^2\alpha^2 = n^2\beta^2$
\[
\omega_m \longmapsto \alpha^2 c'_m(\omega_m) = \alpha^2 m^2\, \frac{c_m(\omega_m)}{12} = \delta^2\, \frac{c_m(\omega_m)}{12} \in \Gamma\, \mathcal{T}_X^{\otimes \delta^2}
\]
\[
\omega'_n \longmapsto \beta^2 c'_n(\omega'_n) = \beta^2 n^2\, \frac{c_n(\omega'_n)}{12} = \delta^2\, \frac{c_n(\omega'_n)}{12} \in \Gamma\, \mathcal{T}_X^{\otimes \delta^2}
\]
\[
G'_{m,n}(\omega_m, \omega'_n) = \alpha^2 c'_m(\omega_m) - \beta^2 c'_n(\omega'_n) \in \Gamma\, \mathcal{T}_X^{\otimes \delta^2}
\]
\[
\begin{cases}
G(t, as) = G(t,s) & a \in \Gamma(S, \underline{\mathcal{O}}_S^*) \\
G(t, s+b) = G(t,s) & b \in \Gamma(S, \underline{\mathcal{O}}_S) \\
G(t, \tfrac{1}{s}) = G(t,s) & s \in \underline{\mathrm{Et}}_S(X, \mathbb{G}_{m,S})
\end{cases}
\]
\note{la fin de la page, des calculs dispersés sur $b_p$, $f_p$, $\frac{b}{b_p}$, $\frac{f}{f_p}$ avec des flèches $f_0 \leftarrow f$, $g + f_0 \leftarrow f$, n'a pas pu être lue de façon suivie}

\page{80}
\note{feuillet de brouillon, traversé de deux longs traits obliques (biffure probable) ; les notations dispersées sont transcrites dans l'ordre de lecture, la disposition est comprimée. Le lien avec les pages précédentes n'est pas explicite}

$D_x \subset T_{X,x} \supset D'_x$, $D'_x \simeq \bar{D}_x$

\begin{tikzcd}[column sep=small, row sep=small, nodes={font=\scriptsize}]
& \bar{X}_x & & \bar{X}_x \\
& U(x) \arrow[u, "\mathbb{C}\text{-an}"] \arrow[d, "\wr"] & B(x) \arrow[l, hook'] \arrow[r] \arrow[d, "\wr"] & V(x) \arrow[u] \arrow[d, "\wr"] \\
& D_x & \hat{T}_x \arrow[l, hook'] \arrow[r] \arrow[d] & D''_x \\
& & \Sigma_x &
\end{tikzcd}

\note{les trois flèches verticales du bas sont marquées « $\mathbb{R}$-anal. » ; au-dessus de $B(x)$ est écrit $\Sigma_\Gamma(x)$ ; en haut de la page, une première version du même diagramme : $D_x \leftarrow U(x) \hookrightarrow B(x) \hookleftarrow V(x) \Rightarrow D'_x$ (biffé), avec $U(x) \to X_x$, $V(x) \to \bar{X}_x$, $B(x)$ au-dessus de $\Sigma(x)$ et $\hat{T}_x$ ($\mathbb{R}$-an), et les mentions $\mathbb{R}$-an, $\mathbb{C}$-an au-dessus des flèches}

$\widetilde{X}_{(x)} \times_{(X, \Pi(x))} \bar{X}_{(x)} =$ \ill{}

$X(x)$, $\Pi_1(X, x)$ \uncertain{ses éléments}, $\widetilde{X}$, $\Pi_x \times \bar{\Pi}_x$ \ill{}

$B(x)_y =$ \uncertain{classes d'équiv.\ de couples} $(c, \xi)$, $\xi \in \hat{T}_{X,x}$, $c \in \Pi(x,y)$, $(c\lambda, \xi) \simeq (c, \lambda\xi)$ ; $(\lambda \in \Pi(x))$

$\hat{T}_X + \mathcal{O}_X$, $\mathcal{V}_X + \underline{\mathcal{O}}_X = \Omega^{-1} \oplus \Omega^0$

$GP(1, \mathbb{C})$

$\check{E} \otimes E = (\Omega^1 + \Omega^0) \otimes (\Omega^{-1} \oplus \Omega^0) = \begin{pmatrix} \Omega^0 & \Omega^{-1} \\ \Omega^1 & \Omega^0 \end{pmatrix}$

$\check{E} \otimes E \otimes \Omega^1 \simeq \begin{pmatrix} \Omega^1 & \Omega^0 \\ \Omega^2 & \Omega^1 \end{pmatrix}$, $\mathcal{V} \otimes \Omega^1 \simeq \begin{pmatrix} \Omega^1 & \Omega^0 \\ \Omega^2 & -\Omega^1 \end{pmatrix}$

\struck{\ill{}} ; $x^2 = 1$, $x = \pm 1$, $g+1$, $3g-3$, \struck{$2g$}, $1$, $4g-2$

$f, g$, $f + \frac{1}{f}$, $\Omega^1 = \omega^{\otimes 2}$, $\frac{f^2+1}{f}$ sans zéros, $g = -\int df^{-1}$

et $\frac{g}{g'} = \varphi$, $\frac{g'}{g} = \frac{1}{\varphi}$, $\log g = \int \frac{1}{\varphi} + \dots$, $g = c \exp \int \frac{1}{\varphi}$

$f \longmapsto \frac{1}{f}$, $f + \frac{1}{f} = \frac{f^2+1}{f} = h$ sans zéros, $f^2 - fh + 1 =$

\end{document}
